\SourceLeaf{313}
\[
\begin{aligned}
\log b={}&\log\beta+\log m\\
&-\tfrac1{24}mm(2\beta\beta-3\tau\tau)\\
&+\tfrac1{2880}m^4(26\beta^4-60\beta\beta\tau\tau+45\tau^4)\\
&-\tfrac1{362880}m^6(502\beta^6-1638\beta^4\tau\tau+1890\beta\beta\tau^4-945\tau^6)\\
&+\tfrac1{29030400}m^8(7102\beta^8-30120\beta^6\tau\tau+49140\beta^4\tau^4\\
&\qquad\qquad-37800\beta\beta\tau^6+14175\tau^8)\\
&-\text{u.\,s.\,w.}
\end{aligned}
\]

\indent Durch Substitution der oben gegebenen Werthe von \(\log m,mm,m^4\) u.\,s.\,w. erhält man hieraus die gesuchte Reihe, welche sich übrigens auch aus der erstern nach \(\beta,\lambda,r\) fortschreitenden unmittelbar ableiten lässt, indem man \(rr-\beta\beta-\tau\tau\) für \(\lambda\lambda\) substituirt.

\Article{19}
Für unsern Zweck reicht es hin, die Formeln nur bis zur vierten Ordnung (einschl.) genau aufzustellen, nemlich
\[
\begin{aligned}
\log t={}&\log\tau+\tfrac1{24}(2rr+\tau\tau)+\tfrac1{2880}(14r^4+20rr\tau\tau+11\tau^4)\\
\log l={}&\log\lambda-\tfrac1{24}(rr-\lambda\lambda)-\tfrac1{2880}(r^4+10rr\lambda\lambda-11\lambda^4)\\
\log b={}&\log\beta-\tfrac1{24}(rr-\beta\beta-3\lambda\lambda)\\
&-\tfrac1{2880}(r^4+10rr\beta\beta+30rr\lambda\lambda-11\beta^4-30\beta\beta\lambda\lambda-45\lambda^4)
\end{aligned}
\]
Anstatt der letzten Formel kann man auch eine der folgenden gebrauchen:
\[
\begin{aligned}
\log b={}&\log\beta+\tfrac1{24}(2rr-2\beta\beta+3\tau\tau)\\
&+\tfrac1{2880}(14r^4-40rr\beta\beta+60rr\tau\tau+26\beta^4-60\beta\beta\tau\tau+45\tau^4)\\
\log b={}&\log\beta+\tfrac1{24}(2\lambda\lambda+\tau\tau)\\
&-\tfrac1{2880}(12\beta\beta\lambda\lambda-12\beta\beta\tau\tau-14\lambda^4-32\lambda\lambda\tau\tau+\tau^4)\\
\log b={}&\log\beta+\tfrac1{24}(2\lambda\lambda+\tau\tau)\\
&-\tfrac1{2880}(12rr\lambda\lambda-12rr\tau\tau-26\lambda^4-8\lambda\lambda\tau\tau-11\tau^4)
\end{aligned}
\]
In allen diesen Formeln sind \(r,\beta,\lambda,\tau,b,l,t\) als in Theilen des Halbmessers ausgedrückt und die Logarithmen als hyperbolische zu verstehen. Sollen dagegen jene sieben Grössen in Bogensecunden ausgedrückt und die Logarithmen die briggischen sein, so erleiden die Formeln weiter keine Veränderung, als dass der gemeinschaftliche Zahlencoëfficient der Glieder zweiter Ordnung \(\frac1{24}\) in \(\frac12\mu\), und der gemeinschaftliche Zahlencoëfficient der Glieder vierter Ordnung \(\frac1{2880}\) in \(\nu\) verwandelt werden muss, wo \(\mu,\nu\) die Producte der Grössen \(\frac1{12}\rho\rho\) und \(\frac1{2880}\rho^4\)
